% S06-M1 (package abhijit/s06-math): commutative algebra. Labels: ss:s06-ca, lem:s06-S, lem:s06-G, lem:s06-D, thm:s06-E,
% if:s06-lemD, ex:s06-toy, def:s06-gor, rem:s06-gor-uses.
\section{Commutative algebra: Gorenstein algebras, congruence ideals, new/old duality}\label{ss:s06-ca}

Throughout this section $A$ is a commutative ring, $S$ a commutative $A$-algebra, and $S^\vee:=\Hom_A(S,A)$. All the results
are pure module theory; the arithmetic of the later sections enters only through their hypotheses. They are the ones the Lean
track formalises in the namespace \pathtt{Langlandscape.CommAlg} (\texttt{CommAlg/\{Gorenstein,LemmaS,NewOld,ThmE,DVR,LemmaG,
BaseChange\}.lean}); we give the Lean name next to each statement (all declarations \stformalized{}, sorry-free, standard axioms only; qualifiers where Lean proves less than the text). Hypotheses of this section enter Lean as explicit arguments; none of the \texttt{Hyp.*} of the later sections is used here.

\subsection{Gorenstein algebras and congruence ideals}

\begin{definition}[traces, Gorenstein, congruence ideal; seed \S2.6]\label{def:s06-gor}
\leavevmode
\begin{enumerate}[label=(\alph*)]
\item The dual $S^\vee$ is an $S$-module by $(s\varphi)(x):=\varphi(sx)$. An element $\ell\in S^\vee$ is a \emph{trace} if $S$ is finite
  free over $A$ and the $S$-linear map $S\to S^\vee$, $s\mapsto s\ell=\ell(s\,\cdot)$, is bijective. Then $S^\vee$ is free of rank one over
  $S$ with generator $\ell$. $S$ is \emph{$A$-Gorenstein} if it has a trace.
  (Lean: \texttt{IsGorensteinWith A S \ensuremath{\ell}} $:=$ \texttt{Module.Free A S \ensuremath{\wedge} Module.Finite A S \ensuremath{\wedge} Function.Bijective
  (fun s => \ensuremath{\ell} \ensuremath{\circ}\ensuremath{_l} LinearMap.mulLeft A s)}, and \texttt{IsGorenstein A S} $:=\exists\ell$, \texttt{IsGorensteinWith A S \ensuremath{\ell}}.)
\item For an ideal $I\subseteq S$ put $\Ann_S(I):=\{s\in S:sI=0\}$ (Lean: \texttt{ann I} $:=(\bot:\mathrm{Ideal}\,S).\mathrm{colon}\,I$).
\item For an $A$-algebra map $\mu:S\to A$ put $S[\mu]:=\Ann_S(\ker\mu)$ and $\eta(S,\mu):=\mu(S[\mu])\subseteq A$, the
  \emph{congruence ideal} (Lean: \texttt{congruenceIdeal \ensuremath{\mu}} $:=(\texttt{ann}\,(\ker\mu)).\mathrm{map}\,\mu$).
\end{enumerate}
\status{definition}
\end{definition}

Since $\ker\mu$ is generated by the elements $x-\mu(x)\cdot1$, an element $s$ lies in $S[\mu]$ iff $xs=\mu(x)s$ for all $x\in S$: $S[\mu]$ is the
$\mu$-eigenspace of $S$ acting on itself. The definition of ``Gorenstein'' asks for a \emph{free} dual; \emph{cyclic} suffices:

\begin{lemma}[cyclic dual]\label{lem:s06-cyclic}
Let $S$ be finite free over $A$. Then $S$ is $A$-Gorenstein iff $S^\vee$ is a cyclic $S$-module, and then every generator of $S^\vee$ is a
trace. \stproved{} \lean{Langlandscape.CommAlg.isGorenstein\_iff\_cyclic}; the clause ``every generator is a trace'' is \lean{Langlandscape.CommAlg.IsGorensteinWith.of\_surjective}. \stformalized{}.
\end{lemma}
\begin{proof}
If $S^\vee=S\ell$, the $A$-linear map $s\mapsto s\ell$ is a surjection $S\to S^\vee$ between free $A$-modules of the same finite rank
$n$ ($S^\vee$ is free with the dual basis). Composing with an $A$-isomorphism $S^\vee\cong A^n\cong S$ gives a surjective endomorphism
of the finitely generated $A$-module $S$, which is injective (Vasconcelos; this is the Orzech property of commutative rings, Lean
\pathtt{OrzechProperty.bijective\_of\_surjective\_of\_injective}). So $s\mapsto s\ell$ is bijective. The converse is clear.
\end{proof}

\begin{lemma}[base change of congruence ideals; seed Lemma S]\label{lem:s06-S}
Let $S$ be finite free and $A$-Gorenstein with trace $\ell$, and let $\mu:S\to A$ be an $A$-algebra map. Then
$S[\mu]=\Ann_S\ker\mu=A\theta$ is free of rank one, where $\theta$ is the unique element with $\ell(\theta x)=\mu(x)$ for all $x\in S$.
Hence $\eta(S,\mu)=(\mu(\theta))$ is principal. For any ring homomorphism $\varphi:A\to A'$ (not necessarily local), $S\otimes_AA'$ is $A'$-Gorenstein and
$\eta(S\otimes A',\mu\otimes A')=\varphi(\eta(S,\mu))A'$. \stproved{} \lean{Langlandscape.CommAlg.LemmaS.ann\_ker\_eq\_span},
\lean{Langlandscape.CommAlg.LemmaS.trace\_theta\_mul} (with \pathtt{theta\_unique}, \pathtt{smul\_theta\_injective}), \lean{Langlandscape.CommAlg.LemmaS.congruenceIdeal\_eq\_span};
base change \lean{Langlandscape.CommAlg.IsGorenstein.baseChange}, \lean{Langlandscape.CommAlg.LemmaS.congruenceIdeal\_baseChange}
(for an arbitrary $A$-algebra $A'$, i.e.\ $\varphi$ not necessarily local). \stformalized{}, full statement.
\end{lemma}
\begin{proof}
\emph{Existence of $\theta$.} $\mu$ is $A$-linear, so $\mu\in S^\vee$, and since $s\mapsto\ell(s\,\cdot)$ is bijective there is a unique
$\theta\in S$ with $\ell(\theta\,\cdot)=\mu$.

\emph{$\theta\in S[\mu]$.} Let $x\in\ker\mu$. For every $s\in S$, $\ell(\theta x\cdot s)=\mu(xs)=\mu(x)\mu(s)=0$. So $\ell(\theta x\,\cdot)=0$,
and injectivity gives $\theta x=0$.

\emph{$S[\mu]\subseteq A\theta$.} Let $s\in S[\mu]$. For $x\in S$, $\ell(sx)=\ell(\mu(x)s)=\mu(x)\ell(s)$, so
$\ell(s\,\cdot)=\ell(s)\mu=\ell(\ell(s)\theta\,\cdot)$, and injectivity gives $s=\ell(s)\theta$.

\emph{Freeness.} If $a\theta=0$ with $a\in A$ then $a\mu=\ell(a\theta\,\cdot)=0$, so $a=a\mu(1)=0$. Thus $S[\mu]=A\theta\cong A$ and
$\eta(S,\mu)=\mu(A\theta)=(\mu(\theta))$.

\emph{Base change.} Put $S'=A'\otimes_AS$, $\mu'=\mathrm{id}\otimes\mu:S'\to A'$, $\ell'=\mathrm{id}\otimes\ell$. $S'$ is finite free over
$A'$, and since $S$ is finite free, $\Hom_{A'}(S',A')=A'\otimes_AS^\vee$. Under this identification the map $s'\mapsto\ell'(s'\,\cdot)$ is
$\mathrm{id}_{A'}\otimes(s\mapsto\ell(s\,\cdot))$, hence bijective. So $\ell'$ is a trace. The element $\theta':=1\otimes\theta$ satisfies
$\ell'(\theta'x')=\mu'(x')$ (both sides are $A'$-linear in $x'$ and agree on $1\otimes x$). By the first part,
$\eta(S',\mu')=(\mu'(\theta'))=(\varphi(\mu(\theta)))=\varphi(\eta(S,\mu))A'$.
\end{proof}
\noindent The base-change argument does not use that $\varphi$ is local; it uses that $\theta$ exists, i.e.\ (Gor$_S$). Without a
Gorenstein hypothesis $\Ann$ does not commute with non-flat base change, and the formula fails in general.

\begin{lemma}[Gorenstein from one fibre; seed Lemma G]\label{lem:s06-G}
Let $A$ be local, $S$ finite free over $A$ and $\mathfrak p\subseteq\m_A$ an ideal. If $S/\mathfrak pS$ is $A/\mathfrak p$-Gorenstein, then
$S$ is $A$-Gorenstein. \stproved{} \lean{Langlandscape.CommAlg.isGorenstein\_of\_isGorenstein\_quotient}. \stformalized{}, full statement.
\end{lemma}
\begin{proof}
Since $S$ is finite free, $S^\vee\otimes_AA/\mathfrak p=\Hom_{A/\mathfrak p}(S/\mathfrak pS,A/\mathfrak p)=(S/\mathfrak pS)^\vee$, compatibly with the
$S$-module structures. Let $\bar\ell$ be a trace of $S/\mathfrak pS$ and $\ell\in S^\vee$ a lift. The finitely generated $A$-module
$N:=S^\vee/S\ell$ satisfies $N=\mathfrak pN\subseteq\m_AN$, because $S\ell$ maps onto $(S/\mathfrak pS)\bar\ell=(S/\mathfrak pS)^\vee$. By
Nakayama's lemma over the local ring $A$, $N=0$, i.e.\ $S^\vee$ is cyclic. Lemma~\ref{lem:s06-cyclic} concludes.
\end{proof}
\noindent \emph{Use.} Lemma~\ref{lem:s06-G} is auxiliary: the chapter does not apply it. It is kept for the Lean track (formalized in \pathtt{CommAlg/LemmaG.lean}) and records how (Gor$_\T$) at one weight lifts to $\Lam$; in the text this passage is inside \cite{Hs}, \cite{MW86} (Wiles' weight-two statement to the $\Lam$-adic Gorenstein property), which we cite instead.
\noindent The seed assumes $S$ local; this is not needed. The tex argues by Nakayama over the semilocal ring $S$; the argument above uses
Nakayama over $A$ only, since $S^\vee$ is a finite $A$-module.

\subsection{New/old duality}

\begin{lemma}[new/old duality; seed Lemma D, corrected]\label{lem:s06-D}
Let $A$ be a Noetherian local domain and $S$ a finite free $A$-algebra, with ideals $\fa,\fb$ such that $S/\fa$, $S/\fb$ are
$A$-torsion-free, $\fa\cap\fb=0$ and $S_K=(S/\fa)_K\times(S/\fb)_K$, where $K=\mathrm{Frac}\,A$. Then:
\begin{enumerate}[label=(\alph*)]
\item $\Ann_S(\fa)=\fb$ and $\Ann_S(\fb)=\fa$. (No Gorenstein hypothesis is needed.)
\item If $S$ is $A$-Gorenstein with trace $\ell$, then $\beta\mapsto\ell(\beta\,\cdot)$ gives $S$-isomorphisms
  $\fb\cong(S/\fa)^\vee$ and $\fa\cong(S/\fb)^\vee$.
\item If moreover $S/\fb$ is $A$-free, then $\fb$ is a direct summand of $S$ as an $A$-module, hence $A$-free, and
  $s\mapsto\ell(s\,\cdot)|_\fb$ induces $S/\fa\cong\fb^\vee$. In particular $S/\fa$ is $A$-free.
\end{enumerate}
All of this commutes with base change $A\to A'$ that preserves the hypotheses. \stproved{}
(Lean: the interface \texttt{if:s06-lemD} below, \lean{Langlandscape.CommAlg.NewOldData.ann\_a}, \lean{Langlandscape.CommAlg.NewOldData.bEquivDualQuot}, \lean{Langlandscape.CommAlg.NewOldData.quotAEquivDualB}, \dots.)
\stformalized{} for (a)--(c), with these qualifiers: the hypothesis $S_K=(S/\fa)_K\times(S/\fb)_K$ enters Lean in the field-free form \texttt{split} (Lemma~\ref{lem:s06-split}; the passage to the product decomposition is not formalized); the $S$-isomorphisms of (b), (c) are formalized as $A$-linear isomorphisms plus the equivariance lemmas \texttt{\_smul}; the final sentence on base change is \emph{not} formalized.
\end{lemma}

\paragraph{The hypotheses as Lean fields.} The hypotheses of Lemma~\ref{lem:s06-D} are exactly the fields of the structure
\texttt{NewOldData A S} of \pathtt{CommAlg/Gorenstein.lean} (which assumes only \texttt{[CommRing A] [CommRing S] [Algebra A S]}),
together with the instance arguments of \pathtt{CommAlg/NewOld.lean}:
\begin{center}\small
\begin{tabular}{@{}p{0.43\textwidth}p{0.52\textwidth}@{}}
\toprule
hypothesis of Lemma~\ref{lem:s06-D} & Lean \\
\midrule
$A$ Noetherian local domain & \texttt{[IsDomain A] [IsNoetherianRing A] [IsLocalRing A]} \\
ideals $\fa,\fb$ of $S$ & fields \texttt{\ensuremath{\mathfrak a} \ensuremath{\mathfrak b} : Ideal S} \\
$S$ finite free over $A$ & fields \texttt{free : Module.Free A S}, \texttt{finite : Module.Finite A S} \\
$S/\fa$, $S/\fb$ torsion-free & \texttt{tf\_a : Module.IsTorsionFree A (S \ensuremath{/} \ensuremath{\mathfrak a})}, \texttt{tf\_b} \\
$\fa\cap\fb=0$ & \texttt{inf\_eq\_bot : \ensuremath{\mathfrak a} \ensuremath{\sqcap} \ensuremath{\mathfrak b} = \ensuremath{\bot}} \\
$S_K=(S/\fa)_K\times(S/\fb)_K$ & \texttt{split : \ensuremath{\exists} c : A, c \ensuremath{\neq} 0 \ensuremath{\wedge} algebraMap A S c \ensuremath{\in} \ensuremath{\mathfrak a} \ensuremath{\sqcup} \ensuremath{\mathfrak b}} \\
(b): $S$ Gorenstein with trace $\ell$ & \texttt{\{\ensuremath{\ell} : Module.Dual A S\} (h\ensuremath{\ell} : IsGorensteinWith A S \ensuremath{\ell})} \\
(c): $S/\fb$ is $A$-free & \texttt{(hb : Module.Free A (S \ensuremath{/} D.\ensuremath{\mathfrak b}))} \\
\bottomrule
\end{tabular}
\end{center}
The field \texttt{split} is the field-free form of the splitting:

\begin{lemma}[the field-free splitting]\label{lem:s06-split}
Let $A$ be a domain with fraction field $K$ and $\fa,\fb\subseteq S$ ideals with $\fa\cap\fb=0$. Then
$S_K\to(S/\fa)_K\times(S/\fb)_K$ is an isomorphism iff $(\fa+\fb)S_K=S_K$ iff there is $c\in A\smallsetminus\{0\}$ with
$c\cdot1\in\fa+\fb$. \stproved{} \lean{Langlandscape.CommAlg.NewOldData.split\_iff}, \stformalized{} \emph{for the second equivalence only} ($\exists c\ne0$, $c\cdot1\in\fa+\fb$ $\iff$ $(\fa+\fb)S_K=S_K$, with $S_K=\mathrm{Frac}A\otimes_AS$); the first equivalence (Chinese remainder theorem, product decomposition) is not formalized.
\end{lemma}
\begin{proof}
Localisation $S\mapsto S_K=K\otimes_AS$ is exact, so $\fa_K\cap\fb_K=(\fa\cap\fb)_K=0$, and $(S/\fa)_K=S_K/\fa_K$. By the Chinese remainder
theorem $S_K\to S_K/\fa_K\times S_K/\fb_K$ is injective, and it is surjective iff $\fa_K+\fb_K=S_K$. Finally $1\in\fa_K+\fb_K=(\fa+\fb)_K$
iff $c\cdot1\in\fa+\fb$ for some non-zero $c\in A$ (clear denominators).
\end{proof}

\begin{proof}[Proof of Lemma~\ref{lem:s06-D}]
Fix $c\ne0$ in $A$ and $x\in\fa$, $y\in\fb$ with $c\cdot1=x+y$ (Lemma~\ref{lem:s06-split}).

(a) $\fb\fa\subseteq\fa\cap\fb=0$, so $\fb\subseteq\Ann_S(\fa)$. Conversely let $s\in\Ann_S(\fa)$. Then $cs=xs+ys=ys\in\fb$, so $c\cdot\bar s=0$ in
$S/\fb$. Since $c\ne0$ and $S/\fb$ is torsion-free, $\bar s=0$, i.e.\ $s\in\fb$. The second equality is symmetric (it uses $S/\fa$
torsion-free). Only \texttt{tf\_a}, \texttt{tf\_b}, \texttt{inf\_eq\_bot}, \texttt{split} and ``$A$ is a domain'' are used.

(b) For $\beta\in\fb$ the functional $\Phi_\fb(\beta):\bar s\mapsto\ell(\beta s)$ on $S/\fa$ is well defined, because $\beta\fa=0$ by (a). The map
$\Phi_\fb:\fb\to(S/\fa)^\vee$ is $A$-linear and satisfies $\Phi_\fb(s\beta)=\Phi_\fb(\beta)\circ(\bar s\cdot)$, which is the $S$-linearity
for the $S$-module structure of Definition~\ref{def:s06-gor}(a) on $(S/\fa)^\vee$. \emph{Injective:} if $\ell(\beta s)=0$ for all $s$ then
$\beta=0$, because $\ell$ is a trace. \emph{Surjective:} let $\psi\in(S/\fa)^\vee$. Then $\psi\circ\pi_\fa\in S^\vee$, so
$\psi\circ\pi_\fa=\ell(\beta\,\cdot)$ for some $\beta\in S$. For $x'\in\fa$ and all $s$, $\ell(\beta x's)=\psi(\pi_\fa(x's))=0$, so $\beta x'=0$.
Hence $\beta\in\Ann_S(\fa)=\fb$ by (a), and $\Phi_\fb(\beta)=\psi$. The isomorphism $\fa\cong(S/\fb)^\vee$ is symmetric.

(c) The sequence $0\to\fb\to S\to S/\fb\to0$ of $A$-modules splits because $S/\fb$ is free, so $S=\fb\oplus C$ as $A$-modules. Thus $\fb$ is a
direct summand of a finite free module, hence finitely generated projective, hence free over the local ring $A$. Define
$\Psi:S/\fa\to\fb^\vee$, $\Psi(\bar s)(\beta)=\ell(s\beta)$; it is well defined because $\fa\fb=0$. \emph{Injective:} if $\ell(s\beta)=0$ for all
$\beta\in\fb$, then for all $t\in S$, $\ell(s\beta\cdot t)=\ell(s\cdot\beta t)=0$ since $\beta t\in\fb$; so $s\beta=0$ for all $\beta\in\fb$, i.e.\
$s\in\Ann_S(\fb)=\fa$. \emph{Surjective:} for $\chi\in\fb^\vee$, $\chi\circ\mathrm{pr}_\fb\in S^\vee$ (projection along $C$) equals
$\ell(s\,\cdot)$ for some $s$, and then $\Psi(\bar s)=\chi$. So $S/\fa\cong\fb^\vee$ is the dual of a finite free module, hence free.

\emph{Base change.} If $A\to A'$ is a map to a Noetherian local domain such that $S\otimes A'$ with the extended ideals $\fa A',\fb A'$
again satisfies the hypotheses, then (a)--(c) apply to it, and the isomorphisms are given by the same formulas with
$\ell'=\mathrm{id}\otimes\ell$ (a trace by Lemma~\ref{lem:s06-S}); so they are the base changes of those over $A$.
\end{proof}
\noindent The seed states (c) without saying \emph{why} $S/\fa$ is free. Freeness of the dual $(S/\fa)^\vee\cong\fb$ alone does not
imply it when $\dim A\ge2$. The argument above uses the summand property, i.e.\ freeness of $S/\fb$ (Lemma~\ref{lem:s06-decomp}(b)).
Only $\fb$ needs to be an $A$-summand; $S$ need not be local (stage 05's ``local $S$'' is harmless but unused).

\begin{interface}[\texttt{if:s06-lemD}]\label{if:s06-lemD}
Lemma~\ref{lem:s06-D} (a)--(c): new/old duality for a finite free $A$-algebra $S$ with ideals $\fa,\fb$ as there. (a) $\Ann_S\fa=\fb$,
$\Ann_S\fb=\fa$ (no Gorenstein hypothesis); (b) for $A$-Gorenstein $S$, $\fb\cong(S/\fa)^\vee$ and $\fa\cong(S/\fb)^\vee$; (c) if
moreover $S/\fb$ is $A$-free, $\fb$ is an $A$-summand and $S/\fa\cong\fb^\vee$ is $A$-free. Status \stproved. Imported by stage 05
(Prop.~subquot uses (a), Prop.~Xiso uses (b)).
\provisional{F16: new interface \texttt{if:s06-lemD}; before: Lemma~D was used by stage 05 as a plain citation, without an interface.}

\emph{Literal Lean form} (fixed signature of \pathtt{stage-06.md} \S0, namespace \pathtt{Langlandscape.CommAlg}, file
\pathtt{CommAlg/NewOld.lean}; variables \texttt{[IsDomain A] [IsNoetherianRing A] [IsLocalRing A] (D : NewOldData A S)
\{\ensuremath{\ell} : Module.Dual A S\} (h\ensuremath{\ell} : IsGorensteinWith A S \ensuremath{\ell}) (hb : Module.Free A (S \ensuremath{/} D.\ensuremath{\mathfrak b}))}):
\begin{center}\small
\begin{tabular}{@{}p{0.2\textwidth}p{0.75\textwidth}@{}}
\toprule
Lemma~\ref{lem:s06-D} & Lean declaration \\
\midrule
(a) & \texttt{NewOldData.ann\_a : ann D.\ensuremath{\mathfrak a} = D.\ensuremath{\mathfrak b}}; \texttt{NewOldData.ann\_b : ann D.\ensuremath{\mathfrak b} = D.\ensuremath{\mathfrak a}} \\
(b), $\fb\cong(S/\fa)^\vee$ & \texttt{NewOldData.bEquivDualQuot : D.\ensuremath{\mathfrak b} \ensuremath{\simeq}\ensuremath{_l}[A] Module.Dual A (S \ensuremath{/} D.\ensuremath{\mathfrak a})} (uses \texttt{h\ensuremath{\ell}}) with
  \texttt{bEquivDualQuot\_apply (b) (s) : D.bEquivDualQuot h\ensuremath{\ell} b (Ideal.Quotient.mk D.\ensuremath{\mathfrak a} s) = \ensuremath{\ell} ((b : S) * s)} and
  \texttt{bEquivDualQuot\_smul (s) (b) : D.bEquivDualQuot h\ensuremath{\ell} (s \ensuremath{\bullet} b) = D.bEquivDualQuot h\ensuremath{\ell} b \ensuremath{\circ}\ensuremath{_l} LinearMap.mulLeft A
  (Ideal.Quotient.mk D.\ensuremath{\mathfrak a} s)} \\
(b), $\fa\cong(S/\fb)^\vee$ & \texttt{NewOldData.aEquivDualQuot : D.\ensuremath{\mathfrak a} \ensuremath{\simeq}\ensuremath{_l}[A] Module.Dual A (S \ensuremath{/} D.\ensuremath{\mathfrak b})}, with
  \texttt{\_apply}, \texttt{\_smul} of the same shape \\
(c), summand & \texttt{NewOldData.isComplemented\_b : IsComplemented (D.\ensuremath{\mathfrak b}.restrictScalars A)} (uses \texttt{hb} only) \\
(c), $\fb$ free & \texttt{NewOldData.free\_b : Module.Free A D.\ensuremath{\mathfrak b}} (uses \texttt{hb} only) \\
(c), $S/\fa\cong\fb^\vee$ & \texttt{NewOldData.quotAEquivDualB : (S \ensuremath{/} D.\ensuremath{\mathfrak a}) \ensuremath{\simeq}\ensuremath{_l}[A] Module.Dual A D.\ensuremath{\mathfrak b}} (uses \texttt{h\ensuremath{\ell} hb}) with
  \texttt{quotAEquivDualB\_mk (s) (b) : D.quotAEquivDualB h\ensuremath{\ell} hb (Ideal.Quotient.mk D.\ensuremath{\mathfrak a} s) b = \ensuremath{\ell} (s * (b : S))} \\
(c), $S/\fa$ free & \texttt{NewOldData.free\_quot\_a : Module.Free A (S \ensuremath{/} D.\ensuremath{\mathfrak a})} (uses \texttt{h\ensuremath{\ell} hb}) \\
\bottomrule
\end{tabular}
\end{center}
Since Mathlib has no $S$-module structure on $\Hom_A(-,A)$, the ``$S$-isomorphism'' of (b) is in Lean the $A$-linear isomorphism
\texttt{bEquivDualQuot} together with the equivariance lemma \pathtt{bEquivDualQuot\_smul}; this is exactly the formula
$\Phi_\fb(s\beta)=\Phi_\fb(\beta)\circ(\bar s\cdot)$ of the proof. The formula of \texttt{\_apply} is $\Phi_\fb(\beta)(\bar s)=\ell(\beta s)$,
and that of \pathtt{quotAEquivDualB\_mk} is $\Psi(\bar s)(\beta)=\ell(s\beta)$.
\end{interface}

\subsection{The combinatorial half}

\begin{theorem}[the combinatorial half; seed Thm~E, hypotheses explicit]\label{thm:s06-E}
Let $S,\fa,\fb$ be as in Lemma~\ref{lem:s06-D}. Put $S'=S/\fa$, and let $\lambda:S\to A$ factor through $\lambda':S'\to A$. Assume
\begin{itemize}
\item[(Gor$_S$)] $S$ is $A$-Gorenstein.
\end{itemize}
Let $\theta$ generate $S[\lambda]$ (Lemma~\ref{lem:s06-S}), and let $x\in S'$ be its image. Then:
\begin{enumerate}[label=(\roman*)]
\item $\theta\in\fb$, $x\in S'[\lambda']$, $\eta(S,\lambda)=(\lambda'(x))$, and $\cc_q:=\lambda(\fb)=\{\psi(x):\psi\in S'^\vee\}$ (the
  \emph{order ideal} of $x$).
\item If $S'[\lambda']$ is free of rank one with a generator $\theta'$ that is unimodular in $S'$ (i.e.\ $S'[\lambda']$ is an
  $A$-direct summand of $S'$), then $\eta(S,\lambda)=\cc_q\cdot\eta(S',\lambda')$.
\item This holds in particular under (Gor$_{S'}$): $S'$ is finite free and $A$-Gorenstein. It also holds whenever $A$ is a DVR
  (\emph{no} hypothesis on $S'$). For $A$ normal, the identity therefore holds after localisation at every height-one prime.
\item In general $\cc=\mathrm{im}(\fb\to S')\cong\fb\cong S'^\vee$. So, if $S'$ is $A$-free (e.g.\ $S/\fb$ free, Lemma~\ref{lem:s06-D}(c)), $\cc$ is a principal ideal of $S'$ if and only if
  (Gor$_{S'}$) holds.
\end{enumerate}
With $S=R$, $A=I_F$: under (Gor$_R$) and (Gor$_{R'}$), $\eta_F=\cc_q(F)\cdot\eta^B_F$, i.e.\ $(h_F)=\cc_q(F)$. \stproved{}
(Lean, namespace \pathtt{Langlandscape.CommAlg.ThmE}: (i) \lean{Langlandscape.CommAlg.ThmE.theta\_mem\_b}, \pathtt{mk\_theta\_mem\_ann}, \pathtt{eta\_eq\_span\_lambda'}, \lean{Langlandscape.CommAlg.ThmE.c\_eq\_orderIdeal};
(ii) \lean{Langlandscape.CommAlg.ThmE.eta\_eq\_c\_mul\_of\_unimodular};
(iii) Gorenstein half \lean{Langlandscape.CommAlg.ThmE.unimodular\_of\_isGorenstein}, \lean{Langlandscape.CommAlg.ThmE.eta\_eq\_c\_mul\_of\_isGorenstein}, DVR half \lean{Langlandscape.CommAlg.ThmE.eta\_eq\_c\_mul\_of\_dvr} (\pathtt{CommAlg/DVR.lean}, with \pathtt{exists\_unimodular\_of\_dvr});
(iv) \lean{Langlandscape.CommAlg.ThmE.c\_principal\_iff\_isGorenstein}.)
\stformalized{} for (i), (ii), (iii) (both halves) and (iv) for $S'$ $A$-free, under the hypotheses of Lemma~\ref{lem:s06-D} as the structure \texttt{NewOldData} (same qualifiers as there). \emph{Not formalized:} the localisation clause of (iii) for $A$ normal (height-one primes), the isomorphism $\mathrm{im}(\fb\to S')\cong\fb\cong S'^\vee$ as a standalone statement for general $S'$, and the final specialisation $S=R$, $A=I_F$ (the abstract statement is formalized, its instantiation to the Hecke rings is not).
\end{theorem}
\begin{proof}
Let $\ell$ be a trace of $S$, so that $\ell(\theta s)=\lambda(s)$ for all $s$ (Lemma~\ref{lem:s06-S}).

(i) $\fa\subseteq\ker\lambda$, so $\theta\fa\subseteq\theta\ker\lambda=0$ and $\theta\in\Ann_S(\fa)=\fb$ (Lemma~\ref{lem:s06-D}(a)). For $s\in S$,
$\bar s x=\overline{s\theta}=\overline{\lambda(s)\theta}=\lambda'(\bar s)x$, so $x\in S'[\lambda']$. $\eta(S,\lambda)=(\lambda(\theta))=(\lambda'(x))$
by Lemma~\ref{lem:s06-S}. For $\beta\in\fb$, $\lambda(\beta)=\ell(\theta\beta)=\ell(\beta\theta)=\Phi_\fb(\beta)(x)$, with $\Phi_\fb:\fb\xrightarrow\sim
S'^\vee$ the isomorphism of Lemma~\ref{lem:s06-D}(b). As $\beta$ runs through $\fb$, $\Phi_\fb(\beta)$ runs through all of $S'^\vee$, so
$\lambda(\fb)=\{\psi(x):\psi\in S'^\vee\}$. (Note that $\Phi_\fb(\theta)=\lambda'$, since $\ell(\theta s)=\lambda(s)$; this is the seed's
normalisation ``$\ell(\theta)=1$'', which the argument does not need.)

(ii) Write $x=h\theta'$ with $h\in A$ (possible since $x\in S'[\lambda']=A\theta'$). Then $\lambda'(x)=h\lambda'(\theta')$. Unimodularity
means that $\{\psi(\theta'):\psi\in S'^\vee\}=A$ (equivalently, there is $\psi_0$ with $\psi_0(\theta')=1$; equivalently $A\theta'$ is a
direct summand). Hence $\cc_q=\{\psi(h\theta')\}=hA$, and
$\eta(S,\lambda)=(h\lambda'(\theta'))=(h)\cdot(\lambda'(\theta'))=\cc_q\cdot\eta(S',\lambda')$, because $\eta(S',\lambda')=\lambda'(A\theta')$.

(iii) \emph{Under (Gor$_{S'}$)} with trace $\ell'$, Lemma~\ref{lem:s06-S} gives $S'[\lambda']=A\theta'$ with $\ell'(\theta'y)=\lambda'(y)$; in
particular $\ell'(\theta')=\lambda'(1)=1$, so $\theta'$ is unimodular, and (ii) applies.

\emph{Over a DVR} (no hypothesis on $S'$). $S'$ is finite and torsion-free over the PID $A$, hence free. The submodule $S'[\lambda']$ is
saturated: if $a\ne0$ and $ay\in S'[\lambda']$, then $a(zy)=0$ for $z\in\ker\lambda'$, so $zy=0$ by torsion-freeness and $y\in S'[\lambda']$.
So $S'/S'[\lambda']$ is torsion-free, hence free, and $S'[\lambda']$ is a direct summand. Its rank is one: localisation is exact and
$\ker\lambda'$ is finitely generated, so $S'[\lambda']_K=S'_K[\lambda'_K]$; in $S_K=S'_K\times S''_K$ (with $S''=S/\fb$) one has
$S_K[\lambda_K]=S'_K[\lambda'_K]\times0$, because $\lambda_K$ kills $(0,1)$, and $S_K[\lambda_K]=K\theta$ is one-dimensional by
Lemma~\ref{lem:s06-S} applied to the $K$-Gorenstein algebra $S_K$. So $S'[\lambda']$ is free of rank one and a summand, and (ii) applies.
\emph{Remark.} The rank-one step uses (Gor$_S$): it is what makes $S_K$ Gorenstein. Only $S'$ needs no hypothesis over a DVR
(cf.\ the Lean forms \pathtt{ThmE.exists\_unimodular\_of\_dvr} and \pathtt{ThmE.eta\_eq\_c\_mul\_of\_dvr} in \pathtt{CommAlg/DVR.lean}, where (Gor$_S$) is an explicit argument).

\emph{$A$ normal.} For a height-one prime $\mathfrak p$, $A_{\mathfrak p}$ is a DVR, and $S_{\mathfrak p}$, $\fa_{\mathfrak p}$, $\fb_{\mathfrak p}$ satisfy the
hypotheses of Lemma~\ref{lem:s06-D} and (Gor$_S$) (localisation is flat; Lemma~\ref{lem:s06-S}). All three ideals localise:
$\eta(S,\lambda)_{\mathfrak p}=\eta(S_{\mathfrak p},\lambda_{\mathfrak p})$ by Lemma~\ref{lem:s06-S}, $\lambda(\fb)_{\mathfrak p}=\lambda_{\mathfrak p}(\fb_{\mathfrak p})$
trivially, and $\eta(S',\lambda')_{\mathfrak p}=\eta(S'_{\mathfrak p},\lambda'_{\mathfrak p})$ because $\Ann$ of a finitely generated ideal commutes with
flat base change. The DVR case gives the identity in $A_{\mathfrak p}$.

(iv) $\fb\to S'$ is injective because $\fb\cap\fa=0$, so $\cc\cong\fb$ as $S$-modules, and $\fb\cong S'^\vee$ by Lemma~\ref{lem:s06-D}(b). The
$S$-action on $\cc$ and on $S'^\vee$ factors through $S'$. Hence $\cc$ is principal iff $S'^\vee$ is a cyclic $S'$-module, iff (when $S'$
is finite free, as under (c) of Lemma~\ref{lem:s06-D}) $S'$ is $A$-Gorenstein (Lemma~\ref{lem:s06-cyclic}).

Finally, Lemma~\ref{lem:s06-decomp}(c) shows that $(R,\fa,\fb)$ over $A=I_F$ satisfies the hypotheses of Lemma~\ref{lem:s06-D}, with
$R''=R/\fb$ free. Under (Gor$_R$) and (Gor$_{R'}$), (iii) gives $\eta_F=\eta(R,\lambda)=\cc_q(F)\cdot\eta(R',\lambda')=\cc_q(F)\eta^B_F$, and
$\eta_F,\eta^B_F$ are principal (Lemma~\ref{lem:s06-S}). The quotient $(h_F)$ of Definition~\ref{def:s06-rings} is defined when
$\eta^B_F\ne0$, and then $(h_F)=\cc_q(F)$ because $I_F$ is a domain.
\end{proof}
\noindent\emph{On $\eta^B_F\ne0$.} $\eta^B_F=(\lambda'(\theta'))$ vanishes iff $\theta'^2=0$, i.e.\ iff the factor of $R'_K$ through which
$\lambda'_K$ factors is not a field. Since $F$ is primitive (H1), this factor is the field $K_F=\mathrm{Frac}\,I_F$ by multiplicity one
for primitive $q$-new eigensystems, i.e.\ the decomposition $\T_\m\otimes\mathrm{Frac}\,I=\mathrm{Frac}\,I\oplus B$ with $B$ the complement \stcited{} \cite[(3.1), p.~18]{Hs} $=$ \cite[Cor.~3.7]{Hi88b} (verified on Hsieh p.~18; note that $R'_K$ is a factor of $R_K$); so $\eta^B_F\ne0$ (and $\eta_F\ne0$). This is implicit in the seed and the tex and is recorded as
open point (N3) in \pathtt{chapters/s06/NOTES.md}; nothing above depends on it except the division.

\begin{remark}[where each Gorenstein hypothesis is used]\label{rem:s06-gor-uses}
\leavevmode
\begin{center}\small
\begin{tabular}{@{}p{0.27\textwidth}p{0.68\textwidth}@{}}
\toprule
hypothesis & used in \\
\midrule
fields of \texttt{NewOldData} & Lemma~\ref{lem:s06-D}(a) (torsion-freeness, $\fa\cap\fb=0$, splitting); Thm~\ref{thm:s06-E}(i) via (a) \\
(Gor$_S$) $=$ (Gor$_R$) & existence of $\theta$ and principality of $\eta_F$ (Lemma~\ref{lem:s06-S}); the isomorphism $\fb\cong S'^\vee$
  (Lemma~\ref{lem:s06-D}(b)), hence Thm~\ref{thm:s06-E}(i) and (iv); the rank-one statement in the DVR case of (iii); base change to
  $I_F$, to $\widetilde\cO_P$, to $A_{\mathfrak p}$ \\
(Gor$_{S'}$) $=$ (Gor$_{R'}$) & only (iii)$\Rightarrow$(ii): unimodularity of $\theta'$ via $\ell'(\theta')=1$; principality of $\eta^B_F$ \\
$S/\fb$ free & Lemma~\ref{lem:s06-D}(c): $S'$ free, so that (Gor$_{S'}$) makes sense and (iv) applies \\
\bottomrule
\end{tabular}
\end{center}
(Gor$_R$) is \emph{indispensable}: without it $\eta_F$ need not be principal and $\fb\not\cong R'^\vee$, and nothing relates $\eta_F$ to
$\fb$. (Gor$_{R'}$) can be weakened to the unimodularity condition of (ii), and is not needed at all over a DVR, e.g.\ at fixed
weight (Cor.~\ref{cor:s06-fixed}). The seed's phrase ``in the setting of Lemma D'' hides (Gor$_R$).
\emph{Sources.} (Gor$_R$) is \cite[\S3.3, p.~17--18]{Hs} \lit{hsieh\_2017\_hida-families-and-p-adic-triple-product-l-functions\_arxiv}{\S3.3,
p.~18; \S7.2, p.~68} \stcited{} under (CR), via \cite[Cor.~2, p.~482]{Wi95}
\lit{wiles\_1995\_modular-elliptic-curves-and-fermats-last-theorem}{Thm 2.1 and Cor.~2, p.~482} and \cite[\S9 Prop.~2]{MW86}.
(Gor$_{R'}$) is \texttt{if:s04-Mord} (\cite[Thm~4.5 and proof, p.~34]{Hs}: $\T^D(N,\chi,\mathbb I)_\m$ is Gorenstein) transported by the
JL isomorphism $\T'\cong\T^{\ord}_B$ (\cite[Rem.~7.8, p.~71]{Hs}: $\T^D$ is the $N^-$-new quotient). Both are base-changed to $I_F$ by
Lemma~\ref{lem:s06-S}. Consistency: Thm~\ref{thm:s06-E} gives $\eta_F\subseteq\eta^B_F$, which is the inclusion
$(\eta^B_F)\supseteq(\eta_F)$ of \cite[Rem.~7.8]{Hs} (\texttt{if:s04-etaB}).
\end{remark}

\begin{example}[the toy algebra]\label{ex:s06-toy}
Let $A$ be a DVR with uniformiser $\varpi$ (e.g.\ $\Zp$; the Lean test \pathtt{LanglandscapeTests/S06/Toy.lean} uses $\Z_p$) and
$t\ge1$. Put
\[S=\{(x,y)\in A^2:x\equiv y\bmod\varpi^t\}=A e_1\oplus Ae_2,\qquad e_1=(1,1),\ e_2=(0,\varpi^t),\]
a finite free $A$-algebra of rank two ($S\cong A[X]/(X(X-\varpi^t))$, $X\mapsto e_2$). Let $\lambda=\mathrm{pr}_1:S\to A$ (the
``new'' eigensystem) and $\fa:=\ker\mathrm{pr}_1=0\times\varpi^tA$, $\fb:=\ker\mathrm{pr}_2=\varpi^tA\times0$.
\begin{itemize}
\item \emph{\texttt{NewOldData}.} $S/\fa\cong A$ via $\mathrm{pr}_1$ and $S/\fb\cong A$ via $\mathrm{pr}_2$ (both surjective), so both are
  free, in particular torsion-free; $\fa\cap\fb=0$; $\varpi^t\cdot1=(\varpi^t,0)+(0,\varpi^t)\in\fa+\fb$, so \texttt{split} holds with
  $c=\varpi^t$.
\item \emph{(Gor$_S$).} $\ell(x,y):=(x-y)/\varpi^t$ is an $A$-linear form on $S$. In the basis $e_1,e_2$ and its dual basis,
  $\ell(e_1\,\cdot)=(\ell(e_1),\ell(e_2))=(0,-1)$ and $\ell(e_2\,\cdot)=(\ell(e_2),\ell(e_2^2))=(-1,-\varpi^t)$, a matrix of determinant
  $-1$. So $\ell$ is a trace.
\item \emph{Lemma S.} $\theta=(\varpi^t,0)$: indeed $\ell(\theta\cdot(x,y))=\ell(\varpi^tx,0)=x=\lambda(x,y)$. $S[\lambda]=\Ann(\fa)=\fb=A\theta$,
  and $\eta(S,\lambda)=(\lambda(\theta))=(\varpi^t)$.
\item \emph{Lemma D.} $\Ann(\fa)=\fb$, $\Ann(\fb)=\fa$; $\Phi_\fb(u\theta)=(\bar s\mapsto\ell(u\theta s))=u\lambda'$ identifies $\fb$ with
  $(S/\fa)^\vee=A\lambda'$.
\item \emph{Thm E.} $S'=S/\fa=A$, $\lambda'=\mathrm{id}$, $\eta(S',\lambda')=A$, and $\cc_q=\lambda(\fb)=(\varpi^t)$. So
  $\eta(S,\lambda)=\cc_q\cdot\eta(S',\lambda')=(\varpi^t)$, and (Gor$_{S'}$) holds trivially.
\end{itemize}
\emph{Where the hypotheses are used:} the trace $\ell$ (Gor$_S$) produces $\theta$ and the identification $\fb\cong(S/\fa)^\vee$; the
hypothesis (Gor$_{S'}$) is vacuous here because $S'=A$. The defect $\varpi^t$ is the congruence $x\equiv y$ between the new
eigensystem $\mathrm{pr}_1$ and the old one $\mathrm{pr}_2$. This is the shape of $R$ near a branch with $\cc_q(F)=(\varpi^t)$, and it
is the non-vacuity test S06-L10a. \stproved{}
\end{example}
